% from ega0-18-fr.tex
\oldpage[0\textsubscript{IV-1}]{54}
L'homomorphisme $h$ définit sur tous les idéaux de $F$ (et en particulier sur
$\mathfrak{R}$) une structure de $G$-bimodule. Soit
$\theta:\mathfrak{R}\to G$ un homomorphisme de $G$-bimodules rendant
commutatif le diagramme (18.1.7.1)~; cela implique que $\theta$ est
\emph{injectif} et que $\mathfrak{J}'\cap\theta(\mathfrak{R})=0$~; comme
$\theta(\mathfrak{R})$ est un idéal bilatère de $G$, et que
$h(\theta(\mathfrak{R}))=\mathfrak{R}$, l'homomorphisme composé
$g\circ h:G\to B$ se factorise en
$G\xrightarrow{q}G/\theta(\mathfrak{R})\xrightarrow{f}B$, où $f$ est
surjectif~; en outre l'image $\mathfrak{J}$ de $\mathfrak{J}'$ dans
$E=G/\theta(\mathfrak{R})$ est le noyau de $f$ (celui de $g\circ h$ étant
$\mathfrak{J}'\oplus\theta(\mathfrak{R})$), et la restriction à
$\mathfrak{J}'$ de l'homomorphisme canonique $q:G\to E$ est injective. On
peut donc former le \emph{produit fibré} $G'=E\times_BF$, et comme les deux
$A$-homomorphismes $q:G\to E$ et $h:G\to F$ sont tels que
$f\circ q=g\circ h$, ils définissent un unique $A$-homomorphisme $u:G'\to G$
par la propriété universelle de $G'$ (18.1.2)~; nous allons voir que $u$ est
\emph{bijectif}. Il suffit de le prouver lorsque $u$ est considéré comme
homomorphisme de $A$-bimodules~; mais on notera alors que $u$ est compatible
avec les filtrations finies sur $G'$ et $G$ formées respectivement par $G'$
et $\mathfrak{J}''=\operatorname{Ker}(G'\to F)$, et par $G$ et
$\mathfrak{J}'$~; en outre, on a vu (18.1.3) que
$\operatorname{gr}_0u:F\to F$ est l'identité et que
$\operatorname{gr}_1u:\mathfrak{J}''\to\mathfrak{J}'$ est bijectif, donc $u$
lui-même est bijectif (Bourbaki, \emph{Alg. comm.}, chap.~III, \S~2,
n\textsuperscript{o}~8, cor.~3 du th.~1).
\end{env}

\subsection{Extensions d'un anneau par un bimodule}
\label{subsection:0.18.2-fr}

\begin{env}[18.2.1]
\label{0.18.2.1-fr}
Soient $E$ un $A$-anneau augmenté sur $B$, $f:E\to B$ l'augmentation,
$\mathfrak{J}=\operatorname{Ker}(f)$ l'idéal d'augmentation. Si l'on a
$\mathfrak{J}^2=0$, $\mathfrak{J}$ est non seulement un $E$-bimodule mais
aussi un $B$-\emph{bimodule} puisque $B$ est isomorphe à
$E/\mathfrak{J}$~; de façon précise, tout $b\in B$ est de la forme $f(x)$
avec $x\in E$, et si $z,z'$ sont dans $\mathfrak{J}$ on a
$(x+z)z'=xz'$ (resp. $z'(x+z)=z'x$) de sorte que la valeur de $xz'$ (resp.
$z'x$) ne dépend pas de l'élément $x$ de $f^{-1}(b)$, et peut s'écrire $bz'$
(resp. $z'b$) ce qui définit la structure de $B$-bimodule considérée.
Inversement, si $\mathfrak{J}$ est muni d'une structure de $B$-bimodule telle
que $xz'=f(x)z'$ et $z'x=z'f(x)$ pour $x\in E$ et
$z'\in\mathfrak{J}$, il est clair que $\mathfrak{J}^2=0$.
\end{env}

\begin{env}[18.2.2]
\label{0.18.2.2-fr}
On appelle \emph{$A$-extension d'un $A$-anneau $B$ par un $B$-bimodule $L$}
une \emph{suite exacte} d'homomorphismes de $A$-bimodules
\[
  0\longrightarrow L\xrightarrow{j}E\xrightarrow{f}B\longrightarrow0
\]
où $E$ est un $A$-anneau, $f$ un $A$-homomorphisme d'anneaux et l'on a, pour
$x\in E$ et $z\in L$,
\[
  j(f(x)z)=xj(z),\qquad j(zf(x))=j(z)x
\]
d'où résulte (18.2.1) que $j(L)$ est un idéal bilatère de carré nul de $E$.
Par abus de langage on dira aussi que $E$ est une extension de $B$ par $L$.
On dit que deux $A$-extensions $E$, $E'$ de $B$ par $L$ sont
\emph{$A$-équivalentes} s'il existe un isomorphisme de $A$-anneaux
$u:E\xrightarrow{\sim}E'$ (dit aussi \emph{$A$-équivalence de
$A$-extensions}) rendant commutatif le diagramme
\begin{equation}
\label{0.18.2.2.1-fr}
\tag{18.2.2.1}
  \xymatrix@C=2.1pc@R=1.2pc{
    & & E\ar[dr]\ar[dd]_u\ar@{}[dd]|{\wr} & & \\
    0\ar[r] & L\ar[ur]\ar[dr] & & B\ar[r] & 0\\
    & & E'\ar[ur] & &
  }
\end{equation}
\end{env}

